% Separate post-CP32 candidate. Complete direct proof; independent audits pending.
\subsection{Invariant covectors and irreducibility of the reduced Lang polar}

Work over an algebraically closed field \(k\) of characteristic two.
Let \(W\) have odd dimension \(N\geq3\), fix a basis \(e_i\),
and use its actual Frobenius-twisted basis \(e_a'\) on \(W^{(2)}\).
Let
\[
 \beta:\Lambda^2W\longrightarrow W^{(2)},\qquad
 \beta(e_i,e_j)=\sum_a t_{ij}^a e_a',\qquad
 t_{ij}^a=t_{ji}^a,\quad t_{ii}^a=0.
 \tag{DE-LP1}
\]
For \(s\in W\) put \(M_s=\beta(s,-)\). Coordinates \(y_i\)
belong to \(W^*\); their squares are the coordinates of an element
of \((W^{(2)})^*\). Define the alternating linear pencil and its
coefficientwise square-root pencil by
\[
 B(\lambda)_{ij}=\sum_a t_{ij}^a\lambda_a,\qquad
 A(y)_{ij}=\sum_a (t_{ij}^a)^{1/2}y_a.
 \quad\text{Thus}\quad A(y)^{[2]}=B(y^{[2]}).
 \tag{DE-LP2}
\]
Only the constant coefficients are square-rooted in \(A\); the
variables are not. Perfectness makes these constants unambiguous.
In the fixed bases the actual Lang equation is
\[
       M_s^t y^{[2]}=B(y^{[2]})s=y.
 \tag{DE-LP3}
\]
Alternation implies \(s^ty=0\) for every solution. Under
\((s,y)\mapsto(a s,a^{-1}y)\), equation (DE-LP3) is preserved.
Its quotient on the nonzero section cone is therefore the Lang
scheme of the Euler bundle: the vector represented by \((s,y)\)
is \(y\otimes s\in\operatorname{ann}(s)\otimes ks\).

Assume that \(A\) has generic rank \(N-1\). Divide its maximal
Pfaffian vector by the common polynomial factor to obtain a primitive
homogeneous radical vector \(\sigma(y)\), of degree \(e\), with
\[
 A(y)\sigma(y)=0,\qquad \gcd(\sigma_0,\ldots,\sigma_{N-1})=1.
\]
The reduced polar polynomial is
\[
                    q(y)=\sum_i y_i\sigma_i(y)^2.
 \tag{DE-LP4}
\]
No coefficient identification between \(A\) and \(B\) is being
suppressed. In particular (DE-LP4) is the polar of the coefficient-
square-root pencil, exactly the one required by (DE-LP3).

\begin{lemma}[A factor yields an actual invariant Lang covector]
\label{lem:DE-Lang-polar-factor-covector}
\status{Candidate; complete direct proof supplied for independent audit.}
Suppose that on a nonempty open \(U\subset\mathbf P(W)\) the
Frobenius (DE-LP3) is bijective on the rank-\((N-1)\) Euler bundle.
Every irreducible factor \(f\) of \(q\) supplies a global invariant
covector on its finite etale Lang local system, of the form
\[
                 y\longmapsto L_s(y)=s^t D y,
                 \qquad D\in\operatorname{Mat}_N(k).
 \tag{DE-LP5}
\]
If this invariant covector vanishes, then
\[
                   B(y^{[2]})\nabla f(y)=0.
 \tag{DE-LP6}
\]
\end{lemma}
\begin{proof}
First \(q\ne0\): each summand in (DE-LP4) has its unique odd
exponent in the indicated coordinate, so these parity classes
cannot cancel. Moreover
\(\partial q/\partial y_i=\sigma_i^2\). Primitivity consequently
makes \(q\) squarefree, since a repeated irreducible factor would
divide every \(\sigma_i\).

On \(k[y_0,\ldots,y_{N-1}]\), define
\[
 \delta_i=\sum_j A(y)_{ji}^2\frac{\partial}{\partial y_j}.
 \tag{DE-LP7}
\]
All coefficients are squares, so their derivatives vanish.
The derivations commute, and \(\delta_i^2=0\): the cross second
derivatives cancel in characteristic two and each pure second
derivative is zero. These statements also hold on the fraction
field. From \(A\sigma=0\), one gets \(\delta_iq=0\).

Write \(q=fg\), with \(\gcd(f,g)=1\). The identity
\(g\delta_i f+f\delta_i g=0\) shows that
\(f\mid\delta_i f\). All factors may be taken homogeneous.
Since \(\delta_i\) increases degree by one, there is a linear
form \(L_i(y)\), possibly zero, such that
\[
                        \delta_i f=fL_i(y).
 \tag{DE-LP8}
\]
Write \(L_i(y)=\sum_jD_{ij}y_j\). Adjoin independent variables
\(s_i\) and declare them constant for the derivations. Then
\(\delta_s=\sum_i s_i\delta_i\) still has square zero, and
\(\delta_s f=fL_s\), where \(L_s=s^tDy\). Applying it twice gives
the polynomial identity
\[
                        \delta_sL_s=L_s^2.
 \tag{DE-LP9}
\]
On (DE-LP3), \(\delta_s y=B(y^{[2]})s=y\). Since \(L_s\)
is linear in \(y\), Euler's degree-one identity gives
\(\delta_sL_s=L_s\) there. Thus \(L_s^2=L_s\) on the actual
Lang scheme. It is additive in \(y\), so defines a morphism of
finite etale group local systems to the constant \(\mathbf F_2\).
It descends from the section cone because (DE-LP5) is invariant
under \((s,y)\mapsto(a s,a^{-1}y)\). This is an actual invariant
covector of the stated local system, not only a differential
identity on a parameter stratum.

If that covector is zero, the Lang fixed vectors on a geometric
generic fibre span \(\operatorname{ann}(s)\) over the algebraic
closure of its function field. Here we use precisely the usual
unit Lang-frame realization on the Frobenius-bijective open;
there is no rank-dropping fibre in this step. Hence
\(D^t s\in ks\) for a generic \(s\). A constant linear
endomorphism preserving every line in a dense projective open
is scalar: comparing coefficients in
\(s_i(D^ts)_j-s_j(D^ts)_i=0\) gives the assertion. Write
\(D=aI\).

Now (DE-LP9), as a polynomial identity before imposing any Lang
equations, becomes
\[
 0=a\,s^tB(y^{[2]})s=a^2(s^ty)^2.
\]
The first equality is alternation; the independent variables
force \(a=0\). Consequently all \(L_i\) vanish and
\(\delta_i f=0\) for every \(i\), exactly (DE-LP6).
\end{proof}

\begin{theorem}[No invariant covectors forces polar irreducibility]
\label{thm:DE-Lang-polar-irreducible}
\status{Candidate; complete direct proof supplied for independent audit.}
In the setting above, suppose that the geometric Lang image on
\(U\) has no invariant covector. Then the reduced polynomial
\(q\) in (DE-LP4) is irreducible over \(k\).
This conclusion does not require a codimension-three base locus
for \(\sigma\), or even birationality of its projective map.
\end{theorem}
\begin{proof}
Suppose that \(f\) is a proper nonconstant irreducible factor of
\(q\). Its invariant covector from the lemma is zero by hypothesis.
Since the generic radical of \(B(y^{[2]})=A(y)^{[2]}\) is
spanned by \(\sigma(y)^{[2]}\), (DE-LP6) implies
\[
                         \nabla f=h\,\sigma^{[2]}
 \tag{DE-LP10}
\]
for a rational function \(h\in k(y)\).
In fact \(h\) is polynomial. In a reduced numerator/denominator
expression for \(h\), its denominator must divide every
\(\sigma_i^2\), because every entry of \(\nabla f\) is
polynomial; their greatest common divisor is one.

If \(h\ne0\), homogeneity gives
\[
              \deg h=\deg f-1-2e=\deg f-\deg q<0,
\]
which is impossible for a polynomial. Thus \(df=0\).
Over the perfect field \(k\), a polynomial with all partial
derivatives zero is a square. This contradicts irreducibility of
the nonconstant \(f\). Therefore no proper irreducible factor
exists, as claimed. The linear case \(e=0\) is included.
\end{proof}

\begin{corollary}[The main Lang component in every fixed odd system]
\label{cor:DE-odd-polar-main-component}
\status{Candidate; uses separately proved geometric Cartier inputs.}
For every fixed nonordinary principally polarized abelian surface
in characteristic two and every odd \(d\geq1\), the reduced
polar (DE-LP4) for its \(|3d\Theta|\) Cartier pencil is
geometrically irreducible. The nonzero Lang incidence has exactly
one irreducible component meeting the open
\(\operatorname{rank}B(y^{[2]})=N-1\), and this component
dominates the section space. Here \(N=9d^2\).
\end{corollary}
\begin{proof}
Work geometrically. The generic maximal dual rank is
Corollary~\ref{cor:fixed-ppav-all-odd-maximal-dual-Cartier}.
The actual nonempty Frobenius-bijective Euler open is supplied
by Proposition~\ref{prop:DE-common-Cartier-incidence} and its
outgoing odd-isogeny application. The vanishing of invariant
covectors is Theorem~\ref{thm:nonordinary-ppav-all-covectors}.
These are exactly the hypotheses of the theorem, with the
coefficient-square-root convention (DE-LP2), so \(q\) is
irreducible.

We also verify that its hypersurface meets the actual maximal
rank open of the original pencil; an omitted common Pfaffian
factor is not ignored. Choose a point \([s_0]\) on the nonempty
common wall of Proposition~\ref{prop:odd-cover-common-rank-wall-nonempty},
with left kernel \(k\lambda\). Then
\[
 \operatorname{rank}B(\lambda)=N-1,\quad
 \operatorname{rad}B(\lambda)=ks_0,\quad
                       \lambda(s_0^{[2]})=0.
\]
Take the coordinatewise square root \(y=\lambda^{[1/2]}\).
The last equality gives \(y^ts_0=0\), while
\(\sigma(y)^{[2]}\) spans the same radical. Thus \(q(y)=0\)
and \(B(y^{[2]})\) has rank \(N-1\), as required.

On this maximal-rank open, the image of the alternating matrix
\(B(y^{[2]})\) is the annihilator of
\(\sigma(y)^{[2]}\). Equation (DE-LP3) is therefore soluble
in \(s\) exactly when \(q(y)=0\). Its solution set in \(s\)
is an affine line with direction
\(k\sigma(y)^{[2]}\). Since \(y\ne0\), none of these solutions
is \(s=0\). On each coordinate chart this is a torsor under a
line bundle over the nonempty irreducible maximal-rank open
of the polar hypersurface. It is consequently irreducible.
The homogeneous scaling already specified descends the same
description to the nonzero projective Lang incidence.

Finally the full Lang equations have Jacobian \(-I\) in \(y\)
over the section cone. Their scheme is etale and quasi-finite
over that cone. The nonempty incidence just constructed is an
open of dimension \(N\), so its image in the cone is nonempty
open and hence dominant. The same holds after projectivizing.
Its irreducibility makes it the unique component meeting this
maximal-dual-rank open. This argument does not identify the
section solution \(s\) with the auxiliary radical vector
\(s_0\) used to prove nonemptiness.
\end{proof}

\begin{remark}[What polar irreducibility leaves open]
\label{rem:DE-Lang-polar-irreducibility-scope}
\status{Candidate scope; exceptional components remain an actual gate.}
The theorem removes the need to establish codimension at least
three for the reduced Pfaffian base locus. The proof is direct
and uses the actual absence of invariant Lang covectors, not
birationality alone. It neither asserts that that base locus has
codimension three nor contradicts the retained birational
height-two examples.

There may still be further irreducible Lang components whose
projection in \(y\) lies entirely in the rank-\((N-3)\) or
lower strata of \(B(y^{[2]})\). Their compatibility equations
are not replaced by the single reduced polar equation there.
No exclusion of those components, no connectedness of the whole
nonzero Lang cover, and no new full-monodromy degree is asserted.
In particular the unique main component is not silently declared
to be the entire Lang cover. All new arguments are symbolic direct
proofs; no finite computation or new external classification is
needed or claimed.
\end{remark}
